\rtmtight
\begin{tblr}{width=\textwidth,colspec={|Q[c,m,2.1cm]|X[l,m]|},hlines,vlines,hspan=minimal,rowsep=1.5pt,colsep=3pt,row{1}={bg=headgray,font=\bfseries}}
\SetCell{c,m}\hfil\logo\hfil & \textbf{POLITEKNIK NEGERI MALANG}\par\textbf{JURUSAN TEKNOLOGI INFORMASI}\par\textbf{PROGRAM STUDI: D4 TEKNIK INFORMATIKA} \\
\SetCell[c=2]{l} \textbf{RENCANA TUGAS MAHASISWA (RTM) DAN RUBRIK PENILAIAN} & \\
Nama Mata Kuliah & Sistem Operasi \\
Kode Mata Kuliah & RTI252004 \quad \textbf{sks / Jam perminggu:} 2 SKS/4 jam \quad \textbf{Semester:} 2 \\
Dosen Pengampu & \begin{enumerate}\item Dr. Yuri Ariyanto, S.Kom., M.Kom.\item Agung Nugroho Pramuditha, S.T., M.T.\item Dian Hanifudin Subhi, S.Kom., M.Kom.\item Muhammad Afif Hendrawan, S.Kom., M.T.\item Imam Fahrur Rozi, S.T., M.T.\end{enumerate} \\
\end{tblr}
\assessmentblock{Tugas Praktikum Instalasi dan Terminal}{Hands-on Practice}{SCPMK0901-01301, SCPMK0206-01302}{Mahasiswa menyelesaikan tugas praktikum instalasi dan terminal sesuai Sub-CPMK dan konteks mata kuliah.}{Menganalisis kebutuhan, mengerjakan artefak praktik/proyek, menguji hasil, mendokumentasikan evidence, dan mempresentasikan keputusan bila diminta.}{Laporan, artefak digital, repository/prototype, evidence pengujian, dan/atau bahan presentasi.}{Kesiapan instalasi dan konfigurasi lingkungan & 8.8 & Instalasi Linux, konfigurasi VM/dual boot, akun, jaringan dasar, dan struktur direktori siap digunakan serta terdokumentasi. & Lingkungan berjalan, tetapi beberapa konfigurasi atau dokumentasi instalasi belum lengkap. & Lingkungan tidak stabil, instalasi gagal, atau tidak dapat digunakan untuk praktikum. \\ \\
Ketepatan penggunaan perintah terminal & 7.7 & Perintah navigasi, file, permission dasar, proses, dan informasi sistem digunakan tepat dengan hasil yang dapat diverifikasi. & Sebagian besar perintah benar, tetapi masih ada kesalahan opsi, urutan, atau interpretasi output. & Perintah terminal banyak keliru atau output tidak dipahami. \\ \\
Evidence praktik dan troubleshooting & 5.5 & Screenshot/log menunjukkan langkah, hasil, masalah, dan solusi troubleshooting secara runtut. & Evidence tersedia, tetapi troubleshooting atau penjelasan hasil masih terbatas. & Evidence tidak lengkap atau tidak membuktikan praktik dilakukan. \\}

\assessmentblock{UTS Administrasi Linux Dasar}{UTS praktik}{SCPMK0901-01301, SCPMK0206-01302}{Mahasiswa menyelesaikan uts administrasi linux dasar sesuai Sub-CPMK dan konteks mata kuliah.}{Menganalisis kebutuhan, mengerjakan artefak praktik/proyek, menguji hasil, mendokumentasikan evidence, dan mempresentasikan keputusan bila diminta.}{Laporan, artefak digital, repository/prototype, evidence pengujian, dan/atau bahan presentasi.}{Penguasaan administrasi dasar Linux & 8 & Instalasi, struktur filesystem, permission, user/group, proses, dan paket dasar diselesaikan tepat pada studi kasus. & Sebagian administrasi dasar benar, tetapi ada konfigurasi atau konsep yang belum konsisten. & Tidak mampu menyelesaikan tugas administrasi dasar secara benar. \\ \\
Ketepatan eksekusi prosedur praktik & 7 & Langkah praktik runtut, command sesuai kebutuhan, dan output diverifikasi tanpa merusak konfigurasi sistem. & Prosedur berjalan, tetapi verifikasi output atau urutan langkah belum lengkap. & Prosedur tidak runtut, banyak command keliru, atau konfigurasi bermasalah. \\ \\
Dokumentasi dan defense hasil administrasi & 5 & Laporan dan defense menjelaskan alasan command, hasil, kendala, dan perbaikan dengan bukti terminal. & Dokumentasi cukup jelas, tetapi alasan command atau defense masih kurang mendalam. & Tidak dapat menjelaskan hasil praktik atau bukti tidak memadai. \\}

\assessmentblock{Tugas Manajemen Akses, Layanan, dan Backup}{Hands-on Practice}{SCPMK0502-01303, SCPMK0208-01304}{Mahasiswa menyelesaikan tugas manajemen akses, layanan, dan backup sesuai Sub-CPMK dan konteks mata kuliah.}{Menganalisis kebutuhan, mengerjakan artefak praktik/proyek, menguji hasil, mendokumentasikan evidence, dan mempresentasikan keputusan bila diminta.}{Laporan, artefak digital, repository/prototype, evidence pengujian, dan/atau bahan presentasi.}{Konfigurasi akses dan keamanan file & 11.2 & User, group, permission, ownership, dan kontrol akses diterapkan sesuai skenario keamanan. & Konfigurasi akses sebagian tepat, tetapi masih ada izin atau ownership yang kurang konsisten. & Konfigurasi akses salah atau berisiko membuka akses yang tidak semestinya. \\ \\
Manajemen layanan dan aplikasi & 9.8 & Service, package, daemon, startup, status, log, dan konfigurasi aplikasi dikelola dengan command yang tepat. & Layanan berjalan, tetapi pengaturan startup, log, atau konfigurasi belum lengkap. & Layanan tidak berjalan atau manajemen aplikasi tidak sesuai skenario. \\ \\
Backup, restore, dan evidence pemulihan & 7 & Backup dibuat, restore diuji, integrity diperiksa, dan bukti pemulihan terdokumentasi jelas. & Backup tersedia, tetapi restore atau pemeriksaan integrity belum lengkap. & Backup tidak dapat dipulihkan atau tidak ada bukti pemulihan. \\}

\assessmentblock{Tugas Besar Remastering Linux}{Project Based Learning}{SCPMK0208-01304}{Mahasiswa menyelesaikan tugas besar remastering linux sesuai Sub-CPMK dan konteks mata kuliah.}{Menganalisis kebutuhan, mengerjakan artefak praktik/proyek, menguji hasil, mendokumentasikan evidence, dan mempresentasikan keputusan bila diminta.}{Laporan, artefak digital, repository/prototype, evidence pengujian, dan/atau bahan presentasi.}{Perencanaan distro remastering & 2 & Tujuan remastering, paket, konfigurasi default, user target, dan kebutuhan sistem dirumuskan jelas. & Rencana remastering cukup jelas, tetapi paket atau konfigurasi belum sepenuhnya sesuai target. & Rencana tidak jelas atau tidak menunjukkan kebutuhan distro hasil remastering. \\ \\
Kualitas image dan konfigurasi hasil remaster & 1.8 & Image dapat boot, paket utama tersedia, konfigurasi desktop/sistem sesuai kebutuhan, dan ukuran terkendali. & Image dapat digunakan, tetapi masih ada paket, konfigurasi, atau konsistensi branding yang kurang. & Image gagal boot atau tidak menunjukkan hasil remastering yang fungsional. \\ \\
Pengujian dan dokumentasi distribusi & 1.2 & Pengujian boot/install, daftar perubahan, kendala, dan panduan penggunaan terdokumentasi lengkap. & Pengujian dilakukan, tetapi cakupan bukti atau panduan penggunaan masih terbatas. & Tidak ada pengujian memadai atau dokumentasi tidak memungkinkan reproduksi. \\}

\assessmentblock{UAS Validasi Administrasi Sistem}{UAS praktik dan defense}{SCPMK0502-01303, SCPMK0208-01304}{Mahasiswa menyelesaikan uas validasi administrasi sistem sesuai Sub-CPMK dan konteks mata kuliah.}{Menganalisis kebutuhan, mengerjakan artefak praktik/proyek, menguji hasil, mendokumentasikan evidence, dan mempresentasikan keputusan bila diminta.}{Laporan, artefak digital, repository/prototype, evidence pengujian, dan/atau bahan presentasi.}{Validasi konfigurasi sistem akhir & 10 & Akses, layanan, aplikasi, backup, dan recovery divalidasi pada skenario administrasi sistem yang realistis. & Validasi mencakup komponen utama, tetapi sebagian skenario atau bukti masih kurang. & Validasi tidak membuktikan kesiapan administrasi sistem. \\ \\
Analisis masalah dan perbaikan & 8.8 & Masalah konfigurasi diidentifikasi dari log/output, diperbaiki tepat, dan dampaknya dijelaskan. & Masalah utama ditemukan, tetapi analisis penyebab atau perbaikan belum lengkap. & Tidak mampu menemukan atau memperbaiki masalah administrasi. \\ \\
Defense praktik administrasi & 6.2 & Keputusan command, konfigurasi, dan trade-off keamanan dijelaskan dengan bukti terminal yang konsisten. & Defense cukup jelas, tetapi beberapa keputusan belum didukung bukti. & Tidak mampu mempertahankan konfigurasi atau bukti praktik tidak valid. \\}

\begin{tblr}{width=\textwidth,colspec={|Q[l,m,3.2cm]|X[l,m]|},hlines,vlines,hspan=minimal,row{1-3}={bg=headgray},rowsep=1.5pt,colsep=3pt}
Jadwal Pelaksanaan: & Minggu 1--17 sesuai RPS. Setiap evaluasi memiliki RTM dan rubrik multi-kriteria. \\
Lain-Lain Yang Diperlukan: & LMS, repositori/prototype/proyek, perangkat lunak sesuai mata kuliah, dan perangkat presentasi. \\
Pustaka: & Mengacu pustaka utama dan pendukung pada bagian RPS. \\
\end{tblr}
